%% LyX 2.1.3 created this file.  For more info, see http://www.lyx.org/.
%% Do not edit unless you really know what you are doing.
%% This specific layout has been created by:
%% Antonio Fernández Ares
%% antares@ugr.es
%% Github @deantares
\documentclass[a4paper,spanish]{scrartcl}
\usepackage[T1]{fontenc}
\usepackage[utf8]{inputenc}
\usepackage{fancyhdr}
\pagestyle{fancy}
\usepackage{color}
\definecolor{note_fontcolor}{rgb}{1, 0.335938, 0}
\usepackage{array}
\usepackage{float}
\usepackage{endnotes}
\usepackage{graphicx}
\usepackage{setspace}
\onehalfspacing

\makeatletter

%%%%%%%%%%%%%%%%%%%%%%%%%%%%%% LyX specific LaTeX commands.
\pdfpageheight\paperheight
\pdfpagewidth\paperwidth

%% Because html converters don't know tabularnewline
\providecommand{\tabularnewline}{\\}
%% The greyedout annotation environment
\newenvironment{lyxgreyedout}
  {\textcolor{note_fontcolor}\bgroup\ignorespaces}
  {\ignorespacesafterend\egroup}

%%%%%%%%%%%%%%%%%%%%%%%%%%%%%% Textclass specific LaTeX commands.
 \let\footnote=\endnote

%%%%%%%%%%%%%%%%%%%%%%%%%%%%%% User specified LaTeX commands.
\usepackage{colortbl}

\definecolor{azulclaro}{rgb}{0.6,0.8,1}

%Aumentar el padding de las tablas
\usepackage{array}
\setlength\extrarowheight{3pt}

%Usar paquete xymatrix
\usepackage[all]{xy}


%Usar paquete para resaltado de código :)
\usepackage{listings}

%Usar paquete para gráficos
\usepackage{tikz}

\usepackage{pgfgantt}

%\newcommand{\paquete}[1]{\[\texttt{#1} \]}

\makeatother

\usepackage{babel}
\addto\shorthandsspanish{\spanishdeactivate{~<>}}

\begin{document}

\lhead{\includegraphics[height=1cm]{imgs/ugrLogo}}


\rhead{\includegraphics[height=1cm]{imgs/posgradoLogo!}}


\title{Plan Investigación}


\author{Paloma de las Cuevas Delgado}


\subtitle{\includegraphics[height=2cm]{imgs/ugrLogo}}

\maketitle
\newpage{}

\begin{center}
\begin{table}[H]
\centering{}%
\begin{tabular}{>{\raggedleft}m{5cm}>{\raggedright}m{8cm}}
\hline 
\multicolumn{2}{|c|}{\textcolor{cyan}{\cellcolor{azulclaro}}Plan de Investigación}\tabularnewline
\hline 
\hline 
\multicolumn{2}{|l|}{\textcolor{cyan}{\cellcolor{azulclaro}} Doctorando:}\tabularnewline
\hline 
Apellidos y nombre: & de las Cuevas Delgado, Paloma \tabularnewline
DNI: & 76418959X\tabularnewline
Correo electrónico: & palomacd@ugr.es\tabularnewline
Firma: & \tabularnewline
 & \tabularnewline
 & \tabularnewline
\hline 
Título provisional de la tesis: & Soft computing techniques applied to corporate and personal security\tabularnewline
Programa de doctorado: & DOCTORADO EN TECNOLOGÍAS DE LA INFORMACIÓN Y LA COMUNICACIÓN (B25.56.1) \tabularnewline
Línea de investigación: & Ingeniería Neuronal y Sistemas Integrados Bioinspirados\tabularnewline
\hline 
\multicolumn{2}{|l}{\textcolor{cyan}{\cellcolor{azulclaro}}Tutor/a:}\tabularnewline
\hline 
Apellidos y nombre: & Merelo Guervós, Juan Julián\tabularnewline
Correo electrónico: & jmerelo@ugr.es\tabularnewline
\hline 
\multicolumn{2}{|l}{\textcolor{cyan}{\cellcolor{azulclaro}}Director/a:}\tabularnewline
\hline 
Apellidos y nombre: & Merelo Guervós, Juan Julián\tabularnewline
Correo electrónico: & jmerelo@ugr.es\tabularnewline
\hline 
Apellidos y nombre: & García Sánchez, Pablo\tabularnewline
Correo electrónico: & pablogarcia@ugr.es \tabularnewline
\hline 
\end{tabular}
\end{table}

\par\end{center}

\newpage{}


\part{Memoria de Investigación}


\section{Título (provisional) de la tesis}

La tesis tiene por título ''Soft computing techniques applied to corporate and personal security''.


\section{Introducción}


\subsection{Proyecto MUSES}


\section{Hipótesis y justificación}


\section{Objetivos}

\subsection{Primer Objetivo}

Definir qué métricas son más interesantes. Esto permitirá que...

\subsection{Segundo Objetivo}

Comparar qué técnicas de Soft Computing son más adecuadas para un entorno como este.

\subsection{Tercer Objetivo}

Desarrollar una metodología (que se pueda vender a la gente) compuesta de varios pasos:
\begin{enumerate}
  \item Pre-procesamiento de datos
  \item Ejecución, comparativa y análisis
  \item Validación de los resultados
\end{enumerate}


\section{Metodología}



\section{Plan de trabajo}

% Dejo esto aquí comentado para aprender cómo se hacen los diagramas de Gantt en LaTeX
%
% \definecolor{barblue}{RGB}{153,204,254}
% \definecolor{groupblue}{RGB}{51,102,254}
% \definecolor{linkred}{RGB}{165,0,33}
% \renewcommand\sfdefault{phv}
% \renewcommand\mddefault{mc}
% \renewcommand\bfdefault{bc}
% \setganttlinklabel{s-s}{EMPEZAR-PARA-EMPEZAR}
% \setganttlinklabel{f-s}{TERMINAR-PARA-EMPEZAR}
% \setganttlinklabel{f-f}{TERMINAR-PARA-TERMINAR}
% \sffamily
% 
% \begin{ganttchart}[
% canvas/.append style={fill=none, draw=black!5, line width=.75pt},
% hgrid style/.style={draw=black!5, line width=.75pt},
% vgrid={*1{draw=black!5, line width=.75pt}},
% today=7,
% today rule/.style={
% draw=black!64,
% dash pattern=on 3.5pt off 4.5pt,
% line width=1.5pt
% },
% today label font=\small\bfseries,
% title/.style={draw=none, fill=none},
% title label font=\bfseries\footnotesize,
% title label node/.append style={below=7pt},
% include title in canvas=false,
% bar label font=\mdseries\small\color{black!70},
% bar label node/.append style={left=2cm},
% bar/.append style={draw=none, fill=black!63},
% bar incomplete/.append style={fill=barblue},
% bar progress label font=\mdseries\footnotesize\color{black!70},
% group incomplete/.append style={fill=groupblue},
% group left shift=0,
% group right shift=0,
% group height=.5,
% group peaks tip position=0,
% group label node/.append style={left=.6cm},
% group progress label font=\bfseries\small,
% link/.style={-latex, line width=1.5pt, linkred},
% link label font=\scriptsize\bfseries,
% link label node/.append style={below left=-2pt and 0pt}
% ]{1}{36}
% \gantttitle[
% title label node/.append style={below left=7pt and -3pt}
% ]{MESES:\quad1}{1}
% \gantttitlelist{2,...,12}{1} \\
% \ganttgroup[progress=57]{\txt{Dispositivo}}{1}{10} \\
% \ganttbar[
% progress=75,
% name=WBS1A
% ]{\textbf{WBS 1.1} Activity A}{1}{8} \\
% \ganttbar[
% progress=67,
% name=WBS1B
% ]{\textbf{WBS 1.2} Activity B}{1}{3} \\
% \ganttbar[
% progress=50,
% name=WBS1C
% ]{\textbf{WBS 1.3} Activity C}{4}{10} \\
% \ganttbar[
% progress=0,
% name=WBS1D
% ]{\textbf{WBS 1.4} Activity D}{4}{10} \\[grid]
% \ganttgroup[progress=0]{WBS 2 Summary Element 2}{4}{10} \\
% \ganttbar[progress=0]{\textbf{WBS 2.1} Activity E}{4}{5} \\
% \ganttbar[progress=0]{\textbf{WBS 2.2} Activity F}{6}{8} \\
% \ganttbar[progress=0]{\textbf{WBS 2.3} Activity G}{9}{10}
% \ganttlink[link type=s-s]{WBS1A}{WBS1B}
% \ganttlink[link type=f-s]{WBS1B}{WBS1C}
% \ganttlink[
% link type=f-f,
% link label node/.append style=left
% ]{WBS1C}{WBS1D}
% \end{ganttchart}


\section{Medios y financiación}


\part{Anexos}

\bibliographystyle{abbrv}
\bibliography{PlanTesis}

\end{document}
